\section{上确界填补有理数的空洞}
\label{sec:12-Real-Analysis-Support/01-Real-Numbers-and-Completeness/01}

古希腊人发现等腰直角三角形的斜边与直角边不可公度。翻译成今天的语言：没有有理数的平方等于 2。这件事让人不安——几何上明明有那么一条线段，长度却写不成两个整数之比。

换个角度看这个困境。列出平方后小于 2 的正有理数：

\[
1,\ 1.4,\ 1.41,\ 1.414,\ 1.4142,\ \ldots
\]

每一步我们都能在有理数范围内多取一位小数，每一项都更大一点，每一项的平方都更靠近 2 但永远不到 2。这个数列没有终点——在有理数世界里它追着一个不存在的数在跑。

用集合的语言会更清晰。令

\[
A=\{\,q\in\Q:q>0,\ q^2<2\,\}.
\]

\(A\) 有上界：比如 2，或者 1.5，或者 1.42——很多有理数都能压住 \(A\) 的全部元素。可一旦你问我``最小的上界是谁？''，有理数就给不出答案了。任给一个候选的有理数 \(r\)：

\begin{itemize}
\tightlist
\item 若 \(r^2<2\)：稍微调大一点（比如把 \(r\) 换成 \(r+1/n\)，取 \(n\) 足够大），平方仍然小于 2，说明 \(r\) 还不是上界；
\item 若 \(r^2>2\)：稍微调小一点，平方仍然大于 2，说明 \(r\) 不是最小上界，因为更小的数同样是上界。
\end{itemize}

严格写出这个``稍微''，用到 \((r\pm\frac1n)^2 = r^2 \pm \frac{2r}{n}+\frac1{n^2}\)。当 \(n\) 充分大时，修正项可以压到任意小，所以 \((r\pm\frac1n)^2\) 可以任意接近 \(r^2\)。对 \(r^2<2\) 的情形，取充分大的 \(n\) 使 \((r+\frac1n)^2<2\)，于是 \(r+\frac1n\in A\) 比 \(r\) 大；对 \(r^2>2\) 的情形，取充分大的 \(n\) 使 \((r-\frac1n)^2>2\)，于是 \(r-\frac1n\) 仍然是 \(A\) 的上界且比 \(r\) 小。

这个困境不来自 \(\sqrt2\) 的特殊性。任何不能写成有理数的极限点，都是有理数系统的一道裂缝。有理数在极限操作下千疮百孔。

\subsection{上界、最大值与上确界：逐级收紧}

这三层概念一个比一个强，差在``是否属于集合''和``是否最小''上。

设 \(A\subseteq\R\) 非空。称 \(M\) 为 \(A\) 的一个上界，如果

\[
a\le M\qquad\text{对所有 }a\in A.
\]

上界只要求压住所有元素，不要求自己属于 \(A\)，也不要求是压住 \(A\) 的数里最小的那个。一个集合如果有一个上界，就有无穷多个上界——任何更大的数也都是上界。

称 \(m\) 为 \(A\) 的最大值，记作 \(m=\max A\)，需要两件事同时成立：\(m\in A\)，且对所有 \(a\in A\) 有 \(a\le m\)。最大值必须住在集合里。开区间 \((0,1)\) 就没有最大值——无论你挑哪个 \(0.999\ldots\) 里的数，总有一个更靠近 1 的数还在区间里。

这就引出上确界。上确界 \(s=\sup A\) 是全体上界中最小的那个。严格写成两条：

\begin{enumerate}
\tightlist
\item \(s\) 是 \(A\) 的上界：对所有 \(a\in A\)，\(a\le s\)；
\item 任何比 \(s\) 小的数都不再是上界：对任意 \(\varepsilon>0\)，存在 \(a\in A\) 使 \(s-\varepsilon<a\)。
\end{enumerate}

第二条等价于说``\(s\) 可以被 \(A\) 中的元素从下方任意逼近''。上确界不需要属于 \(A\)。

拿 \((0,1)\) 检验：上确界是 1，因为 1 压住了区间里所有的数，而任何小于 1 的数（比如 \(1-\varepsilon\)）与 1 之间一定有区间内的点（比如 \(1-\varepsilon/2\)），所以 \(1-\varepsilon\) 不能是上界。但 1 不在 \((0,1)\) 里，因此最大值不存在。换成闭区间 \([0,1]\)，上确界仍是 1，且 1 就在集合里，所以最大值也是 1。

下界、最小值和下确界 \(\inf A\) 是完全对偶的概念，把所有不等号反过来即可：下确界是全体下界中最大的那个。

\subsection{实数的最小上界性质}

实数的完备性可以用一句话表述：

\begin{quote}
每个非空且有上界的实数集合，都在 \(\R\) 中有上确界。
\end{quote}

两个前提缺一不可。空集 \(\varnothing\) 没有上确界——按定义任何实数都压住空集里的所有元素（因为没有元素需要压），于是``最小上界''没有意义。无上界的集合如 \(\R\) 自身或 \(\N\)，也不存在实数上确界。扩展实数系用 \(\pm\infty\) 来兜底，但那套约定的代数规则与实数完全不同，不能随意混用。

现在用这个公理来追击开篇的困境。考虑实数版本的那个集合：

\[
B=\{\,x\in\R:x>0,\ x^2<2\,\}.
\]

\(B\) 非空（\(1\in B\)）且有上界（比如 2），完备性保证 \(s=\sup B\) 在 \(\R\) 中存在。下面证明这个 \(s\) 就是 \(\sqrt2\)——它的平方恰好等于 2。

排除 \(s^2<2\)：若不然，考察 \((s+\frac1n)^2 = s^2+\frac{2s}{n}+\frac1{n^2}\)。当 \(n\to\infty\) 时后两项趋于零，所以 \((s+\frac1n)^2\to s^2<2\)。取充分大的 \(n\) 使 \((s+\frac1n)^2<2\)，则 \(s+\frac1n\in B\)。但 \(s\) 是 \(B\) 的上界，\(s+\frac1n\) 比上界还大且属于 \(B\)——矛盾。

排除 \(s^2>2\)：此时 \((s-\frac1n)^2 = s^2-\frac{2s}{n}+\frac1{n^2}\to s^2>2\)。取充分大的 \(n\) 使 \((s-\frac1n)^2>2\)，则 \(s-\frac1n\) 仍是 \(B\) 的上界（任何 \(x\in B\) 满足 \(x^2<2<(s-\frac1n)^2\)，结合 \(x>0\) 与 \(s-\frac1n>0\) 得 \(x<s-\frac1n\)）。这给出一个比 \(s\) 更小的上界，与 \(s\) 是最小上界矛盾。

于是唯一的可能性是 \(s^2=2\)。实数系统中 \(\sqrt2\) 的存在由完备性担保，不需要构造它然后验算。

这个论证的美感在于：它没有去``算''\(\sqrt2\) 等于几，而是用完备性声明它存在，再用矛盾法确认它的平方必须是 2。存在性来自公理，具体值来自后续推理。

\subsection{单调有界数列为何收敛：完备性的第一个推论}

若数列 \((x_n)\) 单调递增且有上界，令

\[
s=\sup\{\,x_n:n\in\N\,\}.
\]

\(s\) 的存在由完备性保证。对任意 \(\varepsilon>0\)，由上确界定义的第二条，存在某个 \(N\) 使

\[
x_N > s-\varepsilon.
\]

单调性上场：对所有 \(n\ge N\)，

\[
s-\varepsilon < x_N \le x_n \le s.
\]

于是

\[
|x_n-s| < \varepsilon,
\]

即 \(x_n\to s\)。单调递减且有下界的情形完全对称，极限是下确界。

这个证明的骨架极其简洁：完备性负责提供候选极限 \(s\)，单调性负责把``某个项靠近 \(s\)''升级为``尾部全部靠近 \(s\)''。三条腿各司其职，缺任何一条都会塌：

\begin{itemize}
\tightlist
\item 拿走单调性：\((-1)^n\) 有界（上界 1、下界 \(-1\)），来回振荡，不收敛；
\item 拿走有界性：\(x_n=n\) 单调递增，但上确界不存在（在 \(\R\) 中），发散到无穷；
\item 拿走完备性：在 \(\Q\) 中取一列递增有理数逼近 \(\sqrt2\)，它有有理数上界，单调递增，但在 \(\Q\) 中没有极限——因为``上确界属于空间''这一步在 \(\Q\) 中失败。
\end{itemize}

同一个证明在 \(\R\) 中畅通无阻，在 \(\Q\) 中卡在上确界的存在性上。完备性不是一个抽象的哲学声明，它在每一个论证里就是那一行``令 \(s=\sup\{\ldots\}\)''能否合法写出的问题。

\subsection{Archimedean 性质：连续与离散之间的桥梁}

实数还满足这样一个性质：对任意实数 \(x\)，存在自然数 \(n\) 使得 \(n>x\)。等价地，对任意 \(\varepsilon>0\)，存在 \(n\) 使

\[
\frac1n < \varepsilon.
\]

这意味着无论把容许误差定得多小，总有一个自然数的倒数比它更小。大多数 \(\varepsilon\)--\(N\) 证明的最后一步都长这样：找一个 \(N\)，使得

\[
N > \frac{C}{\varepsilon},
\]

然后由 \(n\ge N\) 推出 \(|x_n-L|<\varepsilon\)。这一步能成立的前提正是 Archimedean 性质——如果存在一个比所有 \(1/n\) 都小的``正无穷小''实数，误差就永远压不到它以下。

Archimedean 性质不是独立公理，它可以从完备性推导出来。假设不然：存在某个实数 \(x\) 使所有 \(n\in\N\) 满足 \(n\le x\)。则 \(\N\) 非空且有上界，由完备性，\(s=\sup\N\) 存在。由上确界定义的第二条，存在 \(n\in\N\) 使 \(n > s-1\)。但 \(n+1\) 也是自然数，且 \(n+1 > s\)，与 \(s\) 是 \(\N\) 的上界矛盾。

这个推导告诉我们，在标准的实数系统中，不存在``无限大''的自然数，也不存在``无限小''的正实数。如果你在非标准分析中见到了无穷小量，那是通过扩张实数系统实现的，与我们这里完备性公理所定义的 \(\R\) 不是同一个对象。

\subsection{上确界是存在性工具，不是计算按钮}

知道 \(\sup A\) 存在，与你能否写出它的闭式表达，是两件独立的事。上确界提供的是存在性担保，而不附带计算方案。这恰好是它强大的原因——你可以用它定义对象、建立不等式、证明极限，而无需先把它算出来。

以下是它反复出现的几个场景：

\begin{itemize}
\tightlist
\item Riemann 积分：上和的下确界与下和的上确界，用上下夹逼定义积分值；
\item 概率论中的 essential supremum：忽略零测集后随机变量``实际上''不会超过的上界；
\item 优化：目标函数可能只在无穷远处趋近最优值，最小值永远取不到，但下确界仍然存在且可被序列逼近；
\item 测度论：外测度通过覆盖的下确界构造，内测度通过包含的上确界构造；
\item 函数空间：sup 范数直接拿上确界做距离——\(\norm{ f}_\infty = \sup_x |f(x)|\)。
\end{itemize}

每次使用上确界时，你都还需要继续干活：证明这个上确界满足你需要的方程、不等式或极限关系。``取 sup''是一个合法的操作，但它不是终点。把上确界声明出来只是第一步，真正的论证在于它如何与集合的其他性质交互。
